\uuid{bmfB}
\exo7id{1200}
\titre{exo7 1200}
\auteur{legall}
\organisation{exo7}
\datecreate{1998-09-01}
\isIndication{false}
\isCorrection{false}
\chapitre{Suite}
\sousChapitre{Convergence}
\cpge{mpsi-mp2i.nombres-reels-suites.e, pcsi.nombres-reels-suites.d, ptsi.nombres-reels-suites.d, mpsi-mp2i.nombres-reels-suites.d, pcsi.nombres-reels-suites.c, ptsi.nombres-reels-suites.c, mpsi-mp2i.nombres-reels-suites.a, pcsi.raisonnement-vocabulaire-ensembliste.c, ptsi.raisonnement-vocabulaire-ensembliste.c, bcpst1.suites-reelles}
\module{Analyse}
\niveau{L1}
\difficulte{}

\contenu{
\texte{
Pour tout $  n \in { \Nn}^*  $ on note $  \displaystyle{ u_n =\sum _{k=1}^n
\frac{1}{ k!}}  $ et $  \displaystyle{ v_n=u_n +\frac{1}{ n!n}}  .$ On rapelle que $  \displaystyle{
e=\lim _{n \rightarrow \infty }u_n }  .$
}
\begin{enumerate}
    \item \question{Montrer que les suites $  (u_n)_{n\geq 1}  $ et $  (v_n)_{n\geq 1}  $
sont adjacentes. En d\'eduire une valeur approch\'ee de $  e   $ \`a $  \displaystyle{\frac{1}{
1000}}  .$}
    \item \question{D\'emontrer que $  e  $ est irrationnel.}
\end{enumerate}
}
